\begin{thebibliography}{99}
\addcontentsline{toc}{chapter}{Bibliography}
\bibitem{ref1} R. R. Selvaraju, M. Cogswell, A. Das, R. Vedantam, D. Parikh, and D. Batra, "Grad-CAM: Visual Explanations From Deep Networks via Gradient-Based Localization," in \emph{Proceedings of the IEEE International Conference on Computer Vision (ICCV)}, 2017.

\bibitem{ref2} A. Chattopadhay, A. Sarkar, P. Howlader, and V. N. Balasubramanian, "Grad-CAM++: Generalized Gradient-Based Visual Explanations for Deep Convolutional Networks," in \emph{2018 IEEE Winter Conference on Applications of Computer Vision (WACV)}, 2018, pp. 839-847.

\bibitem{ref3} M. Sundararajan, A. Taly, and Q. Yan, "Axiomatic Attribution for Deep Networks," in \emph{Proceedings of the 34th International Conference on Machine Learning (ICML)}, PMLR, vol. 70, 2017, pp. 3319-3328.

\bibitem{ref4} S. M. Lundberg and S.-I. Lee, "A Unified Approach to Interpreting Model Predictions," in \emph{Advances in Neural Information Processing Systems}, 2017.

\bibitem{ref5} H. Wang, Z. Wang, M. Du, F. Yang, Z. Zhang, S. Ding, P. Mardziel, and X. Hu, "Score-CAM: Score-Weighted Visual Explanations for Convolutional Neural Networks," in \emph{2020 IEEE/CVF Conference on Computer Vision and Pattern Recognition Workshops (CVPRW)}, 2020, pp. 111-119.

\bibitem{ref6} M. D. Zeiler and R. Fergus, "Visualizing and Understanding Convolutional Networks," in \emph{Computer Vision -- ECCV 2014}, Lecture Notes in Computer Science, vol. 8689, Springer, 2014, pp. 818-833.

\bibitem{ref7} M. B. Muhammad and M. Yeasin, "Eigen-CAM: Class Activation Map using Principal Components," in \emph{2020 International Joint Conference on Neural Networks (IJCNN)}, 2020, pp. 1-7.

\bibitem{ref8} M. T. Ribeiro, S. Singh, and C. Guestrin, "'Why Should I Trust You?': Explaining the Predictions of Any Classifier," in \emph{Proceedings of the 22nd ACM SIGKDD International Conference on Knowledge Discovery and Data Mining (KDD)}, 2016, pp. 1135-1144.

\bibitem{ref9} Meta PyTorch, "Captum: Model Interpretability for PyTorch," official documentation and API reference.

\bibitem{ref10} J. Adebayo, J. Gilmer, M. Muelly, I. Goodfellow, M. Hardt, and B. Kim, "Sanity Checks for Saliency Maps," in \emph{Advances in Neural Information Processing Systems}, 2018.

\bibitem{ref11} C.-K. Yeh, C.-Y. Hsieh, A. S. Suggala, D. I. Inouye, and P. Ravikumar, "On the (In)fidelity and Sensitivity of Explanations," in \emph{Advances in Neural Information Processing Systems}, vol. 32, 2019.

\bibitem{ref12} V. Petsiuk, A. Das, and K. Saenko, "RISE: Randomized Input Sampling for Explanation of Black-box Models," in \emph{British Machine Vision Conference (BMVC)}, 2018.

\bibitem{ref13} W. Samek, A. Binder, G. Montavon, S. Lapuschkin, and K.-R. Muller, "Evaluating the Visualization of What a Deep Neural Network Has Learned," \emph{IEEE Transactions on Neural Networks and Learning Systems}, vol. 28, no. 11, pp. 2660-2673, 2017.

\bibitem{ref14} R. C. Fong and A. Vedaldi, "Interpretable Explanations of Black Boxes by Meaningful Perturbation," in \emph{Proceedings of the IEEE International Conference on Computer Vision (ICCV)}, 2017, pp. 3449-3457.

\bibitem{ref15} A. S. Ross, M. C. Hughes, and F. Doshi-Velez, "Right for the Right Reasons: Training Differentiable Models by Constraining their Explanations," in \emph{Proceedings of the Twenty-Sixth International Joint Conference on Artificial Intelligence (IJCAI-17)}, 2017, pp. 2662-2670.

\bibitem{ref16} N. Arun et al., "Assessing the Trustworthiness of Saliency Maps for Localizing Abnormalities in Medical Imaging," \emph{Radiology: Artificial Intelligence}, vol. 3, no. 6, Art. no. e200267, 2021.

\bibitem{ref17} A. Saporta et al., "Benchmarking saliency methods for chest X-ray interpretation," \emph{Nature Machine Intelligence}, vol. 4, pp. 867-878, 2022.

\bibitem{ref18} Stanford Center for Artificial Intelligence in Medicine and Imaging, "CheXlocalize," official dataset page; associated with A. Saporta et al., "Benchmarking saliency methods for chest X-ray interpretation," \emph{Nature Machine Intelligence}, 2022.

\bibitem{ref19} B. H. M. van der Velden, H. J. Kuijf, K. G. A. Gilhuijs, and M. A. Viergever, "Explainable artificial intelligence (XAI) in deep learning-based medical image analysis," \emph{Medical Image Analysis}, vol. 79, 2022, Art. no. 102470.

\bibitem{ref20} H. Chen, C. Gomez, C.-M. Huang, and M. Unberath, "Explainable medical imaging AI needs human-centered design: guidelines and evidence from a systematic review," \emph{npj Digital Medicine}, vol. 5, 2022, Art. no. 156.

\bibitem{ref21} A. DeGrave, J. D. Janizek, and S.-I. Lee, "AI for radiographic COVID-19 detection selects shortcuts over signal," \emph{Nature Machine Intelligence}, vol. 3, pp. 610-619, 2021.

\bibitem{ref22} J. R. Zech, M. A. Badgeley, M. Liu, A. B. Costa, J. J. Titano, and E. K. Oermann, "Variable generalization performance of a deep learning model to detect pneumonia in chest radiographs: A cross-sectional study," \emph{PLOS Medicine}, vol. 15, no. 11, Art. no. e1002683, 2018.

\bibitem{ref23} L. Oakden-Rayner, J. Dunnmon, G. Carneiro, and C. Re, "Hidden stratification causes clinically meaningful failures in machine learning for medical imaging," in \emph{Proceedings of the ACM Conference on Health, Inference, and Learning (CHIL)}, 2020, pp. 151-159.

\bibitem{ref24} M. Reyes, R. Meier, S. Pereira, C. A. Silva, F.-M. Dahlweid, H. von Tengg-Kobligk, R. M. Summers, and R. Wiest, "On the interpretability of artificial intelligence in radiology: challenges and opportunities," \emph{Radiology: Artificial Intelligence}, vol. 2, no. 3, Art. no. e190043, 2020.

\bibitem{ref25} M. Ghassemi, L. Oakden-Rayner, and A. L. Beam, "The false hope of current approaches to explainable artificial intelligence in health care," \emph{The Lancet Digital Health}, vol. 3, no. 11, pp. e745-e750, 2021.

\bibitem{ref26} R. Geirhos, J.-H. Jacobsen, C. Michaelis, R. Zemel, W. Brendel, M. Bethge, and F. A. Wichmann, "Shortcut learning in deep neural networks," \emph{Nature Machine Intelligence}, vol. 2, pp. 665-673, 2020.

\bibitem{ref27} I. Banerjee et al., "'Shortcuts' Causing Bias in Radiology Artificial Intelligence: Causes, Evaluation, and Mitigation," \emph{Journal of the American College of Radiology}, vol. 20, no. 9, pp. 842-851, 2023.

\bibitem{ref28} J. W. Gichoya et al., "AI recognition of patient race in medical imaging: a modelling study," \emph{The Lancet Digital Health}, vol. 4, no. 6, pp. e406-e414, 2022.

\bibitem{ref29} J. Mongan, L. Moy, and C. E. Kahn Jr., "Checklist for Artificial Intelligence in Medical Imaging (CLAIM): A Guide for Authors and Reviewers," \emph{Radiology: Artificial Intelligence}, vol. 2, no. 2, Art. no. e200029, 2020.

\bibitem{ref30} G. S. Collins et al., "TRIPOD+AI statement: updated guidance for reporting clinical prediction models that use regression or machine learning methods," \emph{BMJ}, vol. 385, Art. no. e078378, 2024.

\bibitem{ref31} V. Sounderajah et al., "A quality assessment tool for artificial intelligence-centered diagnostic test accuracy studies: QUADAS-AI," \emph{Nature Medicine}, vol. 27, no. 10, pp. 1663-1665, 2021.

\bibitem{ref32} B. Vasey et al. and the DECIDE-AI expert group, "Reporting guideline for the early-stage clinical evaluation of decision support systems driven by artificial intelligence: DECIDE-AI," \emph{Nature Medicine}, vol. 28, no. 5, pp. 924-933, 2022.

\bibitem{ref33} Society for Imaging Informatics in Medicine (SIIM), "SIIM-ACR Pneumothorax Segmentation Kaggle Challenge," official challenge information page.

\bibitem{ref34} J. P. Cohen et al., "TorchXRayVision: A library of chest X-ray datasets and models," in \emph{Proceedings of the 5th International Conference on Medical Imaging with Deep Learning}, PMLR, vol. 172, 2022, pp. 231-249.

\bibitem{ref35} X. Wang, Y. Peng, L. Lu, Z. Lu, M. Bagheri, and R. M. Summers, "ChestX-Ray8: Hospital-Scale Chest X-Ray Database and Benchmarks on Weakly-Supervised Classification and Localization of Common Thorax Diseases," in \emph{Proceedings of the IEEE Conference on Computer Vision and Pattern Recognition (CVPR)}, 2017, pp. 3462-3471.

\bibitem{ref36} P. Rajpurkar et al., "Deep learning for chest radiograph diagnosis: A retrospective comparison of the CheXNeXt algorithm to practicing radiologists," \emph{PLOS Medicine}, vol. 15, no. 11, Art. no. e1002686, 2018.

\bibitem{ref37} J. Irvin et al., "CheXpert: A Large Chest Radiograph Dataset with Uncertainty Labels and Expert Comparison," in \emph{Proceedings of the AAAI Conference on Artificial Intelligence}, vol. 33, no. 1, 2019, pp. 590-597.

\bibitem{ref38} A. E. W. Johnson et al., "MIMIC-CXR, a de-identified publicly available database of chest radiographs with free-text reports," \emph{Scientific Data}, vol. 6, Art. no. 317, 2019.

\bibitem{ref39} H. Q. Nguyen et al., "VinDr-CXR: An open dataset of chest X-rays with radiologist's annotations," \emph{Scientific Data}, vol. 9, Art. no. 429, 2022.

\bibitem{ref40} F. Wilcoxon, "Individual Comparisons by Ranking Methods," \emph{Biometrics Bulletin}, vol. 1, no. 6, pp. 80-83, 1945.

\bibitem{ref41} S. Holm, "A Simple Sequentially Rejective Multiple Test Procedure," \emph{Scandinavian Journal of Statistics}, vol. 6, no. 2, pp. 65-70, 1979.

\bibitem{ref42} M. Aickin and H. Gensler, "Adjusting for Multiple Testing When Reporting Research Results: The Bonferroni vs Holm Methods," \emph{American Journal of Public Health}, vol. 86, no. 5, pp. 726-728, 1996.

\bibitem{ref43} J. Demšar, "Statistical Comparisons of Classifiers over Multiple Data Sets," \emph{Journal of Machine Learning Research}, vol. 7, pp. 1-30, 2006.

\bibitem{ref44} A. E. Flanders et al., "Construction of a Machine Learning Dataset through Collaboration: The RSNA 2019 Brain CT Hemorrhage Challenge," \emph{Radiology: Artificial Intelligence}, vol. 2, no. 3, Art. no. e190211, 2020.

\bibitem{ref45} A. Dosovitskiy, L. Beyer, A. Kolesnikov, D. Weissenborn, X. Zhai, T. Unterthiner, M. Dehghani, M. Minderer, G. Heigold, S. Gelly, J. Uszkoreit, and N. Houlsby, "An Image Is Worth 16x16 Words: Transformers for Image Recognition at Scale," in \emph{Proc. International Conference on Learning Representations (ICLR)}, 2021.

\bibitem{ref46} S. Abnar and W. Zuidema, "Quantifying Attention Flow in Transformers," in \emph{Proc. 58th Annual Meeting of the Association for Computational Linguistics (ACL)}, 2020, pp. 4190-4197.

\bibitem{ref47} H. Chefer, S. Gur, and L. Wolf, "Transformer Interpretability Beyond Attention Visualization," in \emph{Proc. IEEE/CVF Conference on Computer Vision and Pattern Recognition (CVPR)}, 2021, pp. 782-791.

\bibitem{ref48} W. Xie, X.-H. Li, C. C. Cao, and N. L. Zhang, "ViT-CX: Causal Explanation of Vision Transformers," in \emph{Proc. 32nd International Joint Conference on Artificial Intelligence (IJCAI)}, 2023.

\bibitem{ref49} J. Gildenblat and contributors, \emph{PyTorch library for CAM methods (pytorch-grad-cam)}, GitHub repository, 2021.
\end{thebibliography}
